\chapter{Компьютер из \nt{GLUT}-нейронов}
\label{sec:glutcomputer}

\section{Правила игры}

В главе~\ref{sec:neurontypes} мы увидели, что подавляющее большинство нейронов мозга --- \nt{GLUT}: они только возбуждают, их веса только положительны. Давайте поставим опыт в уме. Возьмём сколько угодно нейронов одного этого типа и попробуем собрать из них компьютер. Не модель мозга, а именно компьютер: с логикой, сумматором, памятью и тактовым генератором. Мы узнаем, что для этого нужно, чего одним возбуждением сделать нельзя и как это обойти.

Сначала договоримся о модели. Настоящий нейрон сложен, и всю эту сложность мы пока выбросим, оставив только блок с рис.~\ref{fig:blackbox}. Время идёт шагами, $t=0,1,2,\dots$; один шаг --- это время, за которое импульс проходит от одного нейрона к следующему, порядка миллисекунды. В каждый шаг нейрон либо выдаёт импульс, $y=1$, либо молчит, $y=0$. Нейрон $i$ складывает импульсы, пришедшие на предыдущем шаге, с весами и сравнивает сумму с порогом:
\begin{empheq}[box=\keybox]{equation}
y_i(t+1) \;=\; \Theta\Bigl(\sum_j w_{ij}\,y_j(t) \;-\; \theta_i\Bigr), \qquad w_{ij}\ge 0 .
\label{eq:glutneuron}
\end{empheq}
Здесь $\Theta(u)=1$ при $u\ge0$ и $0$ иначе, $\theta_i>0$ --- порог нейрона $i$, $w_{ij}$ --- вес связи от нейрона $j$ к нейрону $i$. Единственное ограничение, которое делает задачу интересной, --- неравенство $w_{ij}\ge0$. Это и есть «только \nt{GLUT}». Веса и пороги будем брать целыми: так легче считать, а общности это не убавляет.

Рисовать нейрон будем кружком, в котором написан порог; входящие стрелки подписаны весами, а если подписи нет, вес равен единице.

\section{Логические элементы}

С чего начинается любой компьютер? С вентилей. Возьмём нейрон с двумя входами $x_1,x_2$ единичного веса. При пороге $\theta=1$ он срабатывает, если пришёл хотя бы один импульс: это \emph{ИЛИ}. При пороге $\theta=2$ --- только если пришли оба: это \emph{И}. А если входов три и порог $2$, нейрон срабатывает, когда импульсов пришло больше, чем не пришло: это \emph{большинство}, или мажоритарный элемент (рис.~\ref{fig:gates}).

\begin{figure}[t]
\centering
\begin{tikzpicture}[>=Stealth,
  nrn/.style={circle,draw,very thick,fill=blue!6,minimum size=0.95cm,inner sep=0pt}]
\foreach \x/\th/\name/\n in {0/1/ИЛИ/2,4.2/2/И/2,8.4/2/{большинство}/3}{
  \node[nrn] (N) at (\x,0) {$\th$};
  \ifnum\n=2
    \draw[->,thick,blue!60!black] (\x-1.6,0.45) node[left,black] {$x_1$} -- (N);
    \draw[->,thick,blue!60!black] (\x-1.6,-0.45) node[left,black] {$x_2$} -- (N);
  \else
    \draw[->,thick,blue!60!black] (\x-1.6,0.6) node[left,black] {$x_1$} -- (N);
    \draw[->,thick,blue!60!black] (\x-1.6,0) node[left,black] {$x_2$} -- (N);
    \draw[->,thick,blue!60!black] (\x-1.6,-0.6) node[left,black] {$x_3$} -- (N);
  \fi
  \draw[->,very thick,red!70!black] (N) -- (\x+1.3,0) node[right,black] {$y$};
  \node[below=0.55cm] at (\x,0) {\small \name};}
\end{tikzpicture}
\caption{Три логических элемента из одного \nt{GLUT}-нейрона. Порог написан внутри кружка, все веса равны единице. ИЛИ: $y=x_1\vee x_2$. И: $y=x_1\wedge x_2$. Большинство: $y=1$, если сработали хотя бы два входа из трёх.}
\label{fig:gates}
\end{figure}

Вообще, нейрон с $n$ входами единичного веса и порогом $k$ --- это элемент «хотя бы $k$ из $n$». При $k=1$ это ИЛИ всех входов, при $k=n$ --- И всех входов. Одного нейрона хватает на то, для чего в обычной логике понадобилось бы дерево из многих вентилей. Это первое, что стоит заметить: пороговый элемент мощнее обычного вентиля.

\section{Чего сделать нельзя}

Теперь попробуем сделать НЕ: нейрон, который срабатывает, когда вход молчит, и молчит, когда вход срабатывает. Попробуйте сами подобрать вес и порог. Не получается, и не случайно. Когда вход молчит, сумма в~\eqref{eq:glutneuron} равна нулю, а порог положителен, так что нейрон молчит. Значит, «молчание на входе» не может вызвать импульс. А может быть, НЕ получится из сети в несколько слоёв? Нет, и это можно доказать сразу для всех сетей.

Назовём один набор входных сигналов \emph{не больше} другого, если всюду, где в первом есть импульс, он есть и во втором. Второй набор получается из первого добавлением импульсов.

\begin{keyblock}
\begin{proposition}[Монотонность]
\label{prop:monotone}
Пусть сеть состоит только из нейронов~\eqref{eq:glutneuron} с $w_{ij}\ge0$ и в начальный момент находится в одном и том же состоянии. Если к входным сигналам добавить импульсы, то ни один импульс ни одного нейрона ни в какой момент времени не исчезнет. Иначе говоря, всё, что вычисляет такая сеть, --- \emph{монотонные} функции входов.
\end{proposition}
\end{keyblock}

Доказательство --- индукция по времени. В момент $t=0$ состояния двух прогонов сети совпадают. Пусть к моменту $t$ во втором прогоне у каждого нейрона импульсов не меньше, чем в первом. Тогда у каждого нейрона сумма $\sum_j w_{ij}y_j(t)$ во втором прогоне не меньше: слагаемые неотрицательны, и во втором прогоне их не меньше. Ступенька $\Theta$ не убывает, значит, и $y_i(t+1)$ во втором прогоне не меньше. Индукция завершена.

Следствия неприятные. НЕ --- немонотонная функция: добавили импульс на вход, пропал импульс на выходе. Исключающее ИЛИ тоже немонотонно: при $x_1=1$ добавление $x_2$ гасит выход. Ни то, ни другое сеть из одних \nt{GLUT}-нейронов вычислить не может. А без НЕ набор И, ИЛИ, большинство не является \emph{функционально полным}: из него нельзя собрать произвольную логическую функцию. Есть и следствие, которое понадобится нам для памяти: \emph{импульс нельзя погасить, добавив импульс}. Погасить что-то в такой сети можно только одним способом --- убрать какой-то другой импульс.

Выглядит как тупик. Но у инженеров есть старый способ из него выйти.

\section{Два провода на бит}

Идея такая: будем передавать каждый бит не по одному проводу, а по двум. По первому идёт импульс, если бит равен единице, по второму --- если бит равен нулю. Всегда активен ровно один провод из пары. Бит $x$ передаётся парой $(x,\bar x)$. Такое кодирование называется \emph{парафазным}, или двухпроводным.

Посмотрите, что стало с НЕ. Чтобы получить $\bar x$, не надо ничего вычислять: достаточно поменять провода местами. НЕ стало бесплатным, и нейронов на него не нужно вовсе. А И и ИЛИ в новом коде вычисляются по правилам де Моргана, каждое двумя нейронами:
\begin{empheq}[box=\keybox]{equation}
\begin{aligned}
x\wedge y \;&\to\; \bigl(\;x\wedge y,\;\; \bar x\vee\bar y\;\bigr),\\
x\vee y \;&\to\; \bigl(\;x\vee y,\;\; \bar x\wedge\bar y\;\bigr).
\end{aligned}
\label{eq:dualrail}
\end{empheq}
Справа все операции монотонны, так что утверждение~\ref{prop:monotone} не нарушено. Монотонность никуда не делась: просто выход $\bar x$ растёт, когда растёт вход $\bar x$.

\begin{keyblock}
\begin{proposition}[Полнота парафазной логики]
\label{prop:dualrail}
Если каждый вход подан парой проводов $(x,\bar x)$, то сеть из \nt{GLUT}-нейронов может вычислить любую логическую функцию, и её ответ тоже будет выдан парой проводов. Цена --- примерно вдвое больше нейронов и проводов.
\end{proposition}
\end{keyblock}

Доказательство следует из~\eqref{eq:dualrail}: И, ИЛИ и НЕ вместе функционально полны, а в парафазном коде все три делаются \nt{GLUT}-нейронами. Например, исключающее ИЛИ $x\oplus y=(x\wedge\bar y)\vee(\bar x\wedge y)$ --- два нейрона И и один нейрон ИЛИ, и ещё три нейрона для провода $\overline{x\oplus y}=(x\wedge y)\vee(\bar x\wedge\bar y)$.

Есть одна тонкость: парафазный код должен приходить на вход уже готовым. Если датчик выдаёт только $x$, получить из него $\bar x$ сеть не может: это было бы НЕ. Значит, каждый датчик должен сам выдавать оба провода. Любопытно, что в зрении так и устроено: сигналы от глаза идут по двум параллельным каналам, один из которых отвечает на усиление света, а другой --- на ослабление.

\section{Сумматор}

Теперь соберём главную деталь арифметического устройства --- полный сумматор. У него три входных бита: $a$, $b$ и перенос $c$ из предыдущего разряда. Выходов два: бит суммы $s$ и перенос $c'$ в следующий разряд. Если $n=a+b+c$ --- число единиц среди входов, то $s=1$ при нечётном $n$, а $c'=1$ при $n\ge2$.

Перенос даётся даром: $c'=1$ при $n\ge2$ --- это в точности элемент большинства с рис.~\ref{fig:gates}. Один нейрон. Его парный провод $\bar c'$ --- большинство из $\bar a,\bar b,\bar c$, ещё один нейрон.

С суммой интереснее. Бит суммы равен единице при $n=1$ и при $n=3$. Разложим это условие на пороговые элементы первого слоя: ИЛИ всех трёх входов (сработало при $n\ge1$), большинство инвертированных входов $\bar a,\bar b,\bar c$ (сработало при $n\le1$) и И всех трёх входов (сработало при $n=3$). Выходной нейрон берёт первые два с весом $1$, третий с весом $2$ и имеет порог $2$:
\begin{empheq}[box=\keybox]{equation}
s \;=\; \Theta\bigl(\,\text{ИЛИ}(a,b,c) + \text{БОЛ}(\bar a,\bar b,\bar c) + 2\,\text{И}(a,b,c) - 2\,\bigr).
\label{eq:sum}
\end{empheq}
Проверим все случаи. При $n=0$ сумма равна $0+1+0=1<2$, выход молчит. При $n=1$ --- $1+1+0=2$, выход срабатывает. При $n=2$ --- $1+0+0=1$, молчит. При $n=3$ --- $1+0+2=3$, срабатывает. Всё верно. Провод $\bar s$ устроен зеркально: в~\eqref{eq:sum} надо поменять местами все провода с чертой и без черты.

\begin{figure}[t]
\centering
\begin{tikzpicture}[>=Stealth,
  nrn/.style={circle,draw,very thick,fill=blue!6,minimum size=0.95cm,inner sep=0pt},
  lab/.style={font=\small}]
\node[nrn] (OR)  at (0,2.4) {$1$};  \node[lab,left=1.1cm] (iOR) at (0,2.4) {$a,b,c$};
\node[nrn] (MAJ) at (0,0.8) {$2$};  \node[lab,left=1.1cm] (iMAJ) at (0,0.8) {$a,b,c$};
\node[nrn] (AND) at (0,-0.8) {$3$}; \node[lab,left=1.1cm] (iAND) at (0,-0.8) {$a,b,c$};
\node[nrn] (NMAJ) at (0,-2.4) {$2$};\node[lab,left=1.1cm] (iNMAJ) at (0,-2.4) {$\bar a,\bar b,\bar c$};
\foreach \n in {OR,MAJ,AND,NMAJ}{\draw[->,thick,blue!60!black] (i\n) -- (\n);}
\foreach \n/\t in {OR/ИЛИ,MAJ/БОЛ,AND/И,NMAJ/БОЛ}{\node[lab,anchor=south,gray!40!black] at (\n.north) {\scriptsize \t};}
\node[nrn,fill=red!8] (S) at (4.2,0.4) {$2$};
\draw[->,thick] (OR) -- (S);
\draw[->,thick] (AND) -- node[below,pos=0.55] {\scriptsize $2$} (S);
\draw[->,thick] (NMAJ) .. controls (2.6,-2.4) and (3.2,-0.6) .. (S);
\draw[->,very thick,red!70!black] (S) -- (5.8,0.4) node[right,black] {$s$};
\draw[->,very thick,red!70!black] (MAJ) -- (1.6,0.8) |- (5.8,1.9) node[right,black] {$c'$};
\draw[->,very thick,red!70!black] (NMAJ) -- (1.2,-2.4) |- (5.8,-2.9) node[right,black] {$\bar c'$};
\end{tikzpicture}
\caption{Полный сумматор из \nt{GLUT}-нейронов, половина для проводов без черты. Порог написан в кружке; вес подписан, если он не равен единице. Перенос $c'$ --- один нейрон большинства. Бит суммы $s$ --- выходной нейрон~\eqref{eq:sum}. Вторая половина схемы, для $\bar s$, --- зеркальная копия: ИЛИ, большинство и И от $\bar a,\bar b,\bar c$ и большинство от $a,b,c$. Уже нарисованное большинство от $\bar a,\bar b,\bar c$ служит обеим половинам.}
\label{fig:fulladder}
\end{figure}

Посчитаем нейроны (рис.~\ref{fig:fulladder}). Первый слой: ИЛИ, большинство и И от прямых проводов и те же три от инвертированных --- шесть нейронов. Второй слой: $s$ и $\bar s$ --- два. Всего \emph{восемь нейронов на разряд}, а переносы $c'$ и $\bar c'$ берутся прямо из первого слоя. Глубина схемы --- два шага: сумма готова через две миллисекунды после того, как пришли входы.

Сумматор чисел из $N$ разрядов --- цепочка таких блоков, в которой перенос из разряда $k$ идёт на вход $c$ разряда $k+1$. Перенос проходит один разряд за один шаг, так что последний бит суммы готов через $N+1$ шагов. Мы проверили это на модели: при сложении $7+1$ в четырёхразрядном сумматоре правильный ответ $8$ появляется на пятом шаге и дальше не меняется. До этого выход показывает неверные промежуточные значения, пока перенос бежит по цепочке, --- ровно как в электронном сумматоре. Читать ответ надо после того, как он установился.

\begin{table}[t]
\centering
\small
\begin{tabular}{lccc}
\toprule
Сумматор & Нейронов & Время установления & Для сравнения: процессор \\
\midrule
4 разряда & 32 & 5~мс & \\
32 разряда & 256 & 33~мс & меньше наносекунды \\
\bottomrule
\end{tabular}
\caption{Сумматор из \nt{GLUT}-нейронов, по восемь нейронов на разряд. Один шаг принят равным миллисекунде. Электронный сумматор быстрее в десятки миллионов раз.}
\label{tab:adder}
\end{table}

\section{Память}

Компьютеру нужна память. Самая простая ячейка памяти из \nt{GLUT}-нейрона очевидна: нейрон с порогом $1$, выход которого заведён на его собственный вход. Один раз получив импульс, он будет срабатывать каждый шаг, вечно. Бит записан.

А как его стереть? Утверждение~\ref{prop:monotone} отвечает сразу: никак. Чтобы нейрон замолчал, импульс надо убрать, а убрать можно только другой импульс, который уже поддерживает работу ячейки. Значит, так и надо сделать: пусть ячейка живёт не сама по себе, а пока её поддерживает внешний сигнал \emph{хранения} $h$. Уберём $h$ --- и ячейка погаснет.

Нужна ещё запись: сигнал разрешения записи $e$ и записываемый бит $d$. Хочется одного нейрона, который при $h=1$ повторяет своё прошлое значение, а при $e=1$ копирует $d$. Попробуйте найти для него веса $w_q,w_h,w_d,w_e$ и порог $\theta$. Должны срабатывать случаи «$q$ и $h$» и «$d$ и $e$», то есть $w_q+w_h\ge\theta$ и $w_d+w_e\ge\theta$. Не должны срабатывать случаи «$q$ и $e$ без $d$» (запись нуля поверх единицы) и «$d$ и $h$ без $e$» (запись без разрешения), то есть $w_q+w_e<\theta$ и $w_d+w_h<\theta$. Сложите первые два неравенства и последние два: слева в обоих случаях стоит одна и та же сумма $w_q+w_h+w_d+w_e$, которая должна быть одновременно не меньше $2\theta$ и меньше $2\theta$. Противоречие. Одним нейроном не обойтись.

Двумя слоями --- легко. Нейрон $H$ с порогом $2$ вычисляет «$q$ и $h$», нейрон $W$ с порогом $2$ --- «$d$ и $e$», а нейрон $q$ с порогом $1$ --- ИЛИ от них. Выход $q$ заведён обратно на вход $H$ (рис.~\ref{fig:latch}). Это три нейрона на провод, шесть на парафазный бит.

\begin{figure}[t]
\centering
\begin{tikzpicture}[>=Stealth,
  nrn/.style={circle,draw,very thick,fill=blue!6,minimum size=0.95cm,inner sep=0pt}]
\node[nrn] (H) at (0,1) {$2$};  \node[anchor=east] at (H.west) {\small $H$\,};
\node[nrn] (W) at (0,-1) {$2$}; \node[below left=0.05cm] at (W.south west) {\small $W$};
\node[nrn,fill=red!8] (Q) at (2.6,0) {$1$};
\draw[->,thick] (H) -- (Q); \draw[->,thick] (W) -- (Q);
\draw[->,thick,blue!60!black] (-2.2,0.6) node[left,black] {$h$ (хранить)} -- (H);
\draw[->,thick,blue!60!black] (-2,-0.65) node[left,black] {$d$ (данные)} -- (W);
\draw[->,thick,blue!60!black] (-2,-1.35) node[left,black] {$e$ (писать)} -- (W);
\draw[->,very thick,red!70!black] (Q) -- (4.3,0) node[right,black] {$q$};
\draw[->,thick,red!70!black] (3.5,0) -- (3.5,2.0) -- (0,2.0) -- (H.north);
\end{tikzpicture}
\caption{Ячейка памяти из трёх \nt{GLUT}-нейронов (один провод парафазной пары; второй провод устроен так же, с $\bar d$ вместо $d$). Пока активен $h$, импульс ходит по петле $q\to H\to q$ и бит хранится. Когда активен $e$, в петлю записывается $d$. Если убрать $h$, петля гаснет.}
\label{fig:latch}
\end{figure}

У этой ячейки есть особенность, которой нет у транзисторной защёлки. Петля $q\to H\to q$ имеет длину два шага, так что импульс, записанный в неё, возвращается через шаг. Если $e$ активен только один шаг, в петле будет импульс через шаг, а не каждый шаг. Петля длины $L$ хранит $L$ независимых битов, по одному на каждую фазу: это не защёлка, а \emph{кольцевой сдвиговый регистр}, память на линии задержки. Нам нужен один бит, поэтому будем держать $e$ активным два шага подряд: тогда записаны обе фазы.

\section{Тактовый генератор}

Последняя деталь --- такт. Соединим $L$ нейронов с порогом $1$ в кольцо: выход каждого на вход следующего, последнего --- на вход первого. Запустим один импульс. Он побежит по кольцу и вернётся к началу через $L$ шагов, и так без конца. Каждый нейрон кольца срабатывает раз в $L$ шагов, и все они сдвинуты друг относительно друга: получился многофазный тактовый генератор. Нужен ещё источник постоянной единицы, чтобы подавать на вход сумматора константы: это нейрон с порогом $1$, заведённый сам на себя, --- та самая вечная ячейка, которую мы не смогли стереть. Здесь её вечность как раз кстати.

Из кольца легко получить сигналы управления для памяти. Сигнал записи $e$ --- ИЛИ от двух соседних нейронов кольца: он активен два шага за такт. Сигнал хранения $h$ --- ИЛИ от всех остальных нейронов кольца: он активен всё остальное время. Заметьте, что $h$ и $e$ --- снова парафазная пара: ровно один из них активен в каждый момент.

\section{Собираем счётчик}

Теперь у нас есть всё, и можно собрать настоящее вычислительное устройство. Пусть это будет четырёхразрядный счётчик: регистр $A$, к которому каждый такт прибавляется единица (рис.~\ref{fig:counter}).

\begin{figure}[t]
\centering
\begin{tikzpicture}[>=Stealth,
  blk/.style={draw,very thick,rounded corners=2pt,fill=blue!6,minimum height=1.0cm,align=center,font=\small}]
\node[blk,minimum width=2.4cm] (A) at (0,0) {регистр $A$\\24 нейрона};
\node[blk,minimum width=2.4cm] (ADD) at (4.2,0) {сумматор $+1$\\32 нейрона};
\node[blk,minimum width=2.4cm] (T) at (8.4,0) {регистр $T$\\24 нейрона};
\node[blk,minimum width=2.2cm,fill=red!6] (R) at (2.1,2.6) {кольцо\\16 нейронов};
\node[blk,minimum width=2.2cm,fill=red!6] (C) at (6.3,2.6) {управление\\4 нейрона};
\node[blk,minimum width=1.6cm,fill=gray!10] (ONE) at (4.2,-3.0) {единица\\1 нейрон};
\draw[->,very thick] (A) -- node[above] {\scriptsize 4 бита} (ADD);
\draw[->,very thick] (ADD) -- node[above] {\scriptsize 4 бита} (T);
\draw[->,very thick] (T.south) -- ++(0,-1.2) -| node[pos=0.12,below] {\scriptsize 4 бита} (A.south);
\draw[->,thick] (ONE) -- (ADD);
\draw[->,thick] (R) -- node[above] {\scriptsize 16 фаз} (C);
\draw[->,thick,red!70!black] (C.south) ++(0.5,0) -- node[right] {\scriptsize $h_T,e_T$} (T.north);
\draw[->,thick,red!70!black] (C.south) ++(-0.5,0) -- ++(0,-0.6) -| node[pos=0.3,above] {\scriptsize $h_A,e_A$} (A.north);
\end{tikzpicture}
\caption{Четырёхразрядный счётчик из $101$ \nt{GLUT}-нейрона. Все шины данных --- парафазные: четыре бита занимают восемь проводов. За один такт кольца, $16$ шагов, сумматор вычисляет $A+1$, результат записывается в регистр $T$, а затем копируется из $T$ в $A$.}
\label{fig:counter}
\end{figure}

Зачем два регистра? Сумматор непрерывно вычисляет $A+1$ из текущего содержимого $A$. Если бы мы записывали результат прямо в $A$, то, едва запись началась, на входе сумматора появилось бы новое значение, и за время записи к нему успела бы прибавиться ещё одна единица, а может быть, и две. Это классическая гонка сигналов. Второй регистр её устраняет: такт делится на две фазы. Сначала, пока $A$ неподвижен и сумматор давно установился, открывается запись в $T$. Потом, пока $T$ неподвижен, открывается запись из $T$ в $A$. В кольце из $16$ нейронов мы отвели под запись в $T$ шаги $8$--$9$, а под запись в $A$ --- шаги $12$--$13$. Сумматору нужно пять шагов, чтобы установиться, так что запас по времени большой.

Второе слагаемое сумматора --- число $0$, а перенос в младший разряд --- единица. В парафазном коде ноль --- это молчащий провод $b$ и активный провод $\bar b$, так что на вход сумматора подаются только молчание и выход нейрона постоянной единицы.

\begin{table}[t]
\centering
\small
\begin{tabular}{lr}
\toprule
Узел & Нейронов \\
\midrule
кольцо тактового генератора & 16 \\
сигналы управления $h_T,e_T,h_A,e_A$ & 4 \\
источник единицы & 1 \\
сумматор: 4 разряда по 8 нейронов & 32 \\
регистры $A$ и $T$: 8 битов по 6 нейронов & 48 \\
\midrule
всего & 101 \\
\bottomrule
\end{tabular}
\caption{Бюджет нейронов четырёхразрядного счётчика.}
\label{tab:counter}
\end{table}

Всего получилось $101$ нейрон (таблица~\ref{tab:counter}). Мы промоделировали эту сеть шаг за шагом по правилу~\eqref{eq:glutneuron}, начав с одного импульса в кольце, включённой единицы и нуля в обоих регистрах. Каждые $16$ шагов регистр $A$ показывает следующее число: $1,2,3,\dots,15,0,1,\dots$ При шаге в миллисекунду это шестьдесят с небольшим отсчётов в секунду. Компьютер из одних \nt{GLUT}-нейронов работает.

\section{Чего не хватает}

Работает --- но хрупко. Вспомните утверждение~\ref{prop:monotone}: в сети из одних \nt{GLUT}-нейронов импульс нельзя погасить, добавив импульс. Теперь представьте, что в кольцо тактового генератора случайно попал лишний импульс. Настоящий нейрон время от времени срабатывает сам по себе, так что это вполне реально. В кольце побегут два импульса вместо одного, такт удвоится, и \emph{ничто в сети не может это исправить}: чтобы убрать лишний импульс, нужен нейрон, который гасит импульс в ответ на импульс, а такого у нас нет. То же с ячейками памяти: случайный импульс в петле хранения превращает ноль в единицу, и единица будет жить, пока не снимут сигнал хранения.

Есть и вторая беда. Мы тщательно следили, чтобы в нашей схеме не было петель с усилением, кроме нужных: кольца, ячеек памяти и источника единицы. Но в сети, где связи возникают без инженера, петли положительной обратной связи неизбежны, и, как мы обсуждали в главе~\ref{sec:neurontypes}, активность в них растёт лавиной. В сети из одних \nt{GLUT}-нейронов её нечем остановить.

Обе беды лечатся одним лекарством --- нейроном, который умеет гасить импульс импульсом. Это \nt{GABA}. С ним НЕ становится одним нейроном, парафазный код перестаёт быть обязательным, ячейку памяти можно сбросить сигналом, а лишний импульс в кольце --- подавить. Выходит, что \nt{GABA} нужен не для того, чтобы вычислять, --- вычислять, как мы убедились, можно и без него, --- а для того, чтобы вычисление было устойчивым.
